\ifdefined\gkver\else\def\gkver{student}\fi
\documentclass[\gkver]{gaokaozhenti}
\newcommand{\ccnum}[1]{{\fontspec{Noto Sans CJK JP}#1}}
\begin{document}

\chapter{2016年全国理综Ⅱ}
%% paperType: 理综物理部分

\begin{enumerate}
\item
%% number: 14
%% paperName: 2016年全国理综Ⅱ第 14 题 6 分
%% typeId: 1
%% type: 单选题
%% chapter: 力物体的平衡
%% point: 动态平衡
%% method: 无
%% score: 6
%% degree: 650
%% duplicateId: 0
%% body:
质量为 $m$ 的物体用轻绳 AB 悬挂于天花板上。用水平向左的力 $F$ 缓慢拉动绳的中点 O，如图所示。用 $T$ 表示绳 OA 段拉力的大小，在 O 点向左移动的过程中\xzanswer{A}

\onepicture[w=2.8cm]{figs/14.png}

\fourchoices[answer=A]
{$F$ 逐渐变大，$T$ 逐渐变大}
{$F$ 逐渐变大，$T$ 逐渐变小}
{$F$ 逐渐变小，$T$ 逐渐变大}
{$F$ 逐渐变小，$T$ 逐渐变小}

%% answer: A
%% memo:
\memoanswer{
	设绳 OA 段与竖直方向的夹角为 $\theta$，对 O 点进行受力分析，列平衡方程得 $F=mg\tan\theta$，$T=\frac{mg}{\cos\theta}$，则随 $\theta$ 的逐渐增大，$F$ 逐渐增大，$T$ 逐渐增大，故 A 正确。

	\onepicture[w=2.4cm]{figs/14_an.jpg}
}

\item
%% number: 15
%% paperName: 2016年全国理综Ⅱ第 15 题 6 分
%% typeId: 1
%% type: 单选题
%% chapter: 电场
%% point: 带电粒子的轨迹类问题
%% method: 无
%% score: 6
%% degree: 550
%% duplicateId: 0
%% body:
如图，P 是固定的点电荷，虚线是以 P 为圆心的两个圆．带电粒子 Q 在 P 的电场中运动，运动轨迹与两圆在同一平面内，a、b、c 为轨迹上的三个点。若 Q 仅受 P 的电场力作用，其在 a、b、c 点的加速度大小分别为 $a_{a}$，$a_{b}$，$a_{c}$，速度大小分别为 $v_{a}$，$v_{b}$，$v_{c}$，则\xzanswer{D}

\onepicture[w=3.0cm]{figs/15.png}

\fourchoices[answer=D]
{$a_{a}>a_{b}>a_{c}$，$v_{a}>v_{c}>v_{b}$}
{$a_{a}>a_{b}>a_{c}$，$v_{b}>v_{c}>v_{a}$}
{$a_{b}>a_{c}>a_{a}$，$v_{b}>v_{c}>v_{a}$}
{$a_{b}>a_{c}>a_{a}$，$v_{a}>v_{c}>v_{b}$}

%% answer: D
%% memo:
\memoanswer{
	粒子只受点电荷的库仑力作用，由牛顿第二定律得 $a=\frac{F}{m}$，$F=\frac{kQq}{r^{2}}$，联立得 $a=\frac{kQq}{mr^{2}}\propto\frac{1}{r^{2}}$，由图知 $r_{a}>r_{c}>r_{b}$，故 $a_{b}>a_{c}>a_{a}$；粒子受力指向轨迹的凹侧，由题意可知两虚线圆为等势线，P、Q 带同种电荷，则带电粒子 Q 的动能与电势能之和不变，因粒子在三点的电势能关系为 $E_{pa}<E_{pc}<E_{pb}$，则有 $E_{ka}>E_{kc}>E_{kb}$，则 $v_{a}>v_{c}>v_{b}$，D 正确。
}

\item
%% number: 16
%% paperName: 2016年全国理综Ⅱ第 16 题 6 分
%% typeId: 1
%% type: 单选题
%% chapter: 曲线运动
%% point: 竖直平面圆周运动
%% method: 无
%% score: 6
%% degree: 550
%% duplicateId: 0
%% body:
小球 P 和 Q 用不可伸长的轻绳悬挂在天花板上，P 球的质量大于 Q 球的质量，悬挂 P 球的绳比悬挂 Q 球的绳短．将两球拉起，使两绳均被水平拉直，如图所示，将两球由静止释放，在各自轨迹的最低点\xzanswer{C}

\onepicture[w=7cm]{figs/16.png}

\fourchoices[answer=C]
{P 球的速度一定大于 Q 球的速度}
{P 球的动能一定小于 Q 球的动能}
{P 球所受绳的拉力一定大于 Q 球所受绳的拉力}
{P 球的向心加速度一定小于 Q 球的向心加速度}

%% answer: C
%% memo:
\memoanswer{
	小球摆动至最低点的过程中，由动能定理得 $mgl=\frac{1}{2}mv^{2}$，得 $v=\sqrt{2gl}$，绳越长速度越大，因 $l_{Q}>l_{P}$，则 Q 球速度大，故 A 错误；结合 A 项分析，动能等于 $mgl$，因为 P 球质量大而绳长短，则无法确定 P、Q 球动能的大小关系，故 B 错误；在最低点，由牛顿第二定律得 $T-mg=m\frac{v^{2}}{l}$，得 $T=3mg$，则质量大的球所受绳的拉力大，故 C 正确；在最低点球的向心加速度 $a=\frac{v^{2}}{l}=2g$，P、Q 球的向心加速度相等，与球的质量和绳长无关，故 D 错误。
}

\item
%% number: 17
%% paperName: 2016年全国理综Ⅱ第 17 题 6 分
%% typeId: 1
%% type: 单选题
%% chapter: 电路
%% point: 电容
%% method: 无
%% score: 6
%% degree: 450
%% duplicateId: 0
%% body:
阻值相等的四个电阻，电容器 $C$ 及电池 $E$（内阻可忽略）连接成如图所示电路．开关 S 断开且电流稳定时，$C$ 所带的电荷量为 $Q_{1}$；闭合开关 S，电流再次稳定后，$C$ 所带的电荷量为 $Q_{2}$。$Q_{1}$ 与 $Q_{2}$ 的比值为\xzanswer{C}

\onepicture[w=5cm]{figs/17.png}

\fourchoices[answer=C]
{$\frac{2}{5}$}
{$\frac{1}{2}$}
{$\frac{3}{5}$}
{$\frac{2}{3}$}

%% answer: C
%% memo:
\memoanswer{
	开关 S 断开与开关 S 闭合，电流稳定时的电路简化图如图所示。

	\onepicture[w=3cm]{figs/17_an1.jpg}
	\onepicture[w=3cm]{figs/17_an2.jpg}

	设每个电阻阻值均为 $R$，开关 S 断开时，并联部分的总电阻为 $\frac{2R\times R}{2R+R}=\frac{2}{3}R$，由闭合电路欧姆定律得并联部分的电压为 $\frac{E}{\frac{2}{3}R+R}\times\frac{2}{3}R=\frac{2}{5}E$，电容器两端的电压等于并联部分电压的一半，故 $U_{C1}=\frac{1}{5}E$，由 $C=\frac{Q}{U}$ 知所带的电荷量 $Q_{1}=\frac{1}{5}CE$；开关 S 闭合时，并联部分的总电阻为 $\frac{R\times R}{R+R}=\frac{1}{2}R$，由闭合电路欧姆定律得并联部分的电压为 $\frac{E}{\frac{1}{2}R+R}\times\frac{1}{2}R=\frac{1}{3}E$，电容器两端的电压等于并联部分电压，故 $U_{C2}=\frac{1}{3}E$，由 $C=\frac{Q}{U}$ 知所带的电荷量 $Q_{2}=\frac{1}{3}CE$，则 $\frac{Q_{1}}{Q_{2}}=\frac{3}{5}$，故 C 正确。
}

\item
%% number: 18
%% paperName: 2016年全国理综Ⅱ第 18 题 6 分
%% typeId: 1
%% type: 单选题
%% chapter: 磁场
%% point: 带电粒子在磁场中的运动
%% method: 无
%% score: 6
%% degree: 450
%% duplicateId: 0
%% body:
一圆筒处于磁感应强度大小为 $B$ 的匀强磁场中，磁场方向与筒的轴平行，筒的横截面如图所示。图中直径 MN 的两端分别开有小孔。筒绕其中心轴以角速度 $\omega$ 顺时针转动．在该截面内，一带电粒子从小孔 M 射入筒内，射入时的运动方向与 MN 成 $30^\circ$ 角。当筒转过 $90^\circ$ 时，该粒子恰好从小孔 N 飞出圆筒。不计重力。若粒子在筒内未与筒壁发生碰撞，则带电粒子的比荷为\xzanswer{A}

\onepicture[w=3.0cm]{figs/18.png}

\fourchoices[answer=A]
{$\frac{\omega}{3B}$}
{$\frac{\omega}{2B}$}
{$\frac{\omega}{B}$}
{$\frac{2\omega}{B}$}

%% answer: A
%% memo:
\memoanswer{
	圆筒转过 $90^\circ$ 所用的时间为 $t=\frac{\frac{\pi}{2}}{\omega}=\frac{\pi}{2\omega}$，小孔 N 顺时针转过 $90^\circ$，带电粒子仍从 N 点射出，由几何关系得带电粒子运动轨迹对应的圆心角为 $\left(\frac{\pi}{3}-\frac{\pi}{4}\right)\times2=30^\circ$，则 $t=\frac{30^\circ}{360^\circ}\cdot\frac{2\pi m}{qB}$，联立可得 $\frac{q}{m}=\frac{\omega}{3B}$，故 A 正确。

	\onepicture[w=3.5cm]{figs/18_an.jpg}
}

\item
%% number: 19
%% paperName: 2016年全国理综Ⅱ第 19 题 6 分
%% typeId: 2
%% type: 多选题
%% chapter: 运动和力
%% point: 变加速直线运动
%% method: 无
%% score: 6
%% degree: 350
%% duplicateId: 0
%% body:
两实心小球甲和乙由同一种材料制成，甲球质量大于乙球质量．两球在空气中由静止下落，假设它们运动时受到的阻力与球的半径成正比，与球的速率无关。若它们下落相同的距离，则\xzanswer{BD}

\fourchoices[answer=BD]
{甲球用的时间比乙球长}
{甲球末速度的大小大于乙球末速度的大小}
{甲球加速度的大小小于乙球加速度的大小}
{甲球克服阻力做的功大于乙球克服阻力做的功}

%% answer: BD
%% memo:
\memoanswer{
	由牛顿第二定律得 $a=\frac{mg-f}{m}=g-\frac{f}{m}$，由题意知 $f=kr$，$m=\rho\frac{4}{3}\pi r^{3}$，联立得 $a=g-\frac{3k}{4\pi\rho r^{2}}$，又因为甲的质量大于乙的质量，则甲的半径大于乙的半径，故甲的加速度大小大于乙的加速度大小，故 C 错误；由匀变速直线运动的规律有 $h=\frac{1}{2}at^{2}$，得 $t=\sqrt{\frac{2h}{a}}$，则 $t_{\text{甲}}<t_{\text{乙}}$，故 A 错误；由匀变速直线运动的规律有 $v^{2}=2ah$，得 $v=\sqrt{2ah}$，则 $v_{\text{甲}}>v_{\text{乙}}$，故 B 正确；克服阻力做功为 $W_{f}=fh=krh$，因甲的半径大于乙的半径，故甲克服阻力做功多，故 D 正确。
}

\item
%% number: 20
%% paperName: 2016年全国理综Ⅱ第 20 题 6 分
%% typeId: 2
%% type: 多选题
%% chapter: 电磁感应
%% point: 切割磁感线电动势
%% method: 无
%% score: 6
%% degree: 350
%% duplicateId: 0
%% body:
法拉第圆盘发动机的示意图如图所示．铜圆盘安装在竖直的铜轴上，两铜片 P、Q 分别与圆盘的边缘和铜轴接触．圆盘处于方向竖直向上的匀强磁场 $B$ 中。圆盘旋转时，关于流过电阻 $R$ 的电流，下列说法正确的是\xzanswer{AB}

\onepicture[w=4.5cm]{figs/20.png}

\fourchoices[answer=AB]
{若圆盘转动的角速度恒定，则电流大小恒定}
{若从上向下看，圆盘顺时针转动，则电流沿 a 到 b 的方向流动}
{若圆盘转动方向不变，角速度大小发生变化，则电流方向可能发生变化}
{若圆盘转动的角速度变为原来的两倍，则电流在 $R$ 上的热功率也变为原来的 2 倍}

%% answer: AB
%% memo:
\memoanswer{
	圆盘切割磁感线产生的感应电动势 $E=Br\frac{\omega r}{2}\propto\omega$，感应电流 $I=\frac{E}{R}\propto\omega$，即圆盘转动的角速度恒定，电流大小恒定，故 A 正确；圆盘切割磁感线，相当于圆盘圆心与 P 点间的半径切割磁感线，由右手定则，电流沿 a 到 b 的方向流动，故 B 正确；由楞次定律知感应电流的方向与圆盘转动的角速度大小无关，故 C 错误；由 A 项分析知 $I\propto\omega$，又 $P=I^{2}R\propto\omega^{2}$，角速度变为原来的两倍，则 R 上的热功率变为原来的 4 倍，故 D 错误。
}

\item
%% number: 21
%% paperName: 2016年全国理综Ⅱ第 21 题 6 分
%% typeId: 2
%% type: 多选题
%% chapter: 功和能
%% point: 动能定理
%% method: 无
%% score: 6
%% degree: 250
%% duplicateId: 0
%% body:
如图，小球套在光滑的竖直杆上，轻弹簧一端固定于 O 点，另一端与小球相连．现将小球从 M 点由静止释放，它在下降的过程中经过了 N 点，已知在 M、N 两点处，弹簧对小球的弹力大小相等。且 $\angle ONM<\angle OMN<\frac{\pi}{2}$，在小球从 M 点运动到 N 点的过程中\xzanswer{BCD}

\onepicture[w=2.2cm]{figs/21.png}

\fourchoices[answer=BCD]
{弹力对小球先做正功后做负功}
{有两个时刻小球的加速度等于重力加速度}
{弹簧长度最短时，弹力对小球做功的功率为零}
{小球到达 N 点时的动能等于其在 M、N 两点的重力势能差}

%% answer: BCD
%% memo:
\memoanswer{
	初态弹簧处于压缩状态，末态弹簧处于伸长状态，且弹簧弹力大小相等，则弹簧弹力先增大再减小再增大，由弹簧弹力与速度方向间的夹角变化可知弹簧弹力对小球先做负功再做正功再做负功，故 A 错误；小球运动过程中受重力、弹簧的弹力、杆的弹力，其中杆的弹力始终垂直于杆，弹簧的弹力沿弹簧方向，当弹簧与光滑杆垂直时，小球竖直方向只受重力的作用，故加速度为重力加速度；当弹簧为原长时，小球只受重力作用，小球的加速度也为重力加速度，故 B 正确；当弹簧与光滑杆垂直时，弹簧长度最短，弹簧弹力与速度垂直，则弹力对小球做功的功率为零，故 C 正确；M、N 两点弹簧弹性势能相等，从 M 到 N 小球的重力势能转化为动能，则小球在 N 点的动能等于其在 M、N 两点的重力势能差，故 D 正确。
}

\item
%% number: 22
%% paperName: 2016年全国理综Ⅱ第 22 题 6 分
%% typeId: 4
%% type: 实验题
%% chapter: 直线运动
%% point: 打点计时器公式
%% method: 无
%% score: 6
%% degree: 550
%% duplicateId: 0
%% body:
某物理小组对轻弹簧的弹性势能进行探究，实验装置如图 \ref{2016全国理综Ⅱ22a} 所示：轻弹簧放置在光滑水平桌面上，弹簧左端固定，右端与一物块接触而不连接，纸带穿过打点计时器并与物块连接。向左推物块使弹簧压缩一段距离，由静止释放物块，通过测量和计算，可求得弹簧被压缩后的弹性势能。
\begin{enumerate}
	\item
	实验中涉及到下列操作步骤：
	\begin{steps}
		\item
		把纸带向左拉直
		\item
		松手释放物块
		\item
		接通打点计时器电源
		\item
		向左推物块使弹簧压缩，并测量弹簧压缩量
	\end{steps}
	上述步骤正确的操作顺序是 \tkanswer[2.5cm]{}（填入代表步骤的序号）。
	\item
	图\ref{2016全国理综Ⅱ22b}中 M 和 L 纸带是分别把弹簧压缩到不同位置后所得到的实际打点结果。打点计时器所用交流电的频率为 $50\UHz$。由 M 纸带所给的数据，可求出在该纸带对应的实验中物块脱离弹簧时的速度为 \tkanswer[2cm]{} $\Ums$。比较两纸带可知，\tkanswer[1.5cm]{}（填“M”或“L”）纸带对应的实验中弹簧被压缩后的弹性势能大。
\end{enumerate}
\twopicture[labelA=2016全国理综Ⅱ22a, labelB=2016全国理综Ⅱ22b, wA=5.5cm, wB=8.5cm, align=right]
{figs/22a.png}
{figs/22b.png}
%% answer: 
\jdanswer{
\begin{enumerate}
	\item
	\ccnum{④①③②}

	\item
	$1.29\Ums$，M

\end{enumerate}
}
%% memo:
\memoanswer{
	（1）先压缩弹簧，松手后，弹簧的弹性势能转化为物块离开弹簧时的动能，测量出动能就可得到弹性势能，保证纸带拉直且先接通电源后释放纸带，则正确步骤是\ccnum{④①③②}。

	（2）物块离开弹簧后，由于桌面光滑，物块做匀速直线运动。打点周期为 $0.02\Us$，取 M 纸带中最后两段求平均值，则 $\overline{v}=\frac{\Delta x}{\Delta t}=\frac{(2.58+2.57)\times10^{-2}}{0.02\times2}\Ums\approx1.29\Ums$，因为纸带 M 匀速段相邻打点间距大，故纸带 M 对应的实验中弹簧的弹性势能大。
}

\item
%% number: 23
%% paperName: 2016年全国理综Ⅱ第 23 题 9 分
%% typeId: 4
%% type: 实验题
%% chapter: 电路
%% point: 电路实验
%% method: 无
%% score: 9
%% degree: 450
%% duplicateId: 0
%% body:
某同学利用图\ref{2016全国理综Ⅱ23a}所示电路测量量程为 $2.5\UV$ 的电压表 V 的内阻（内阻为数千欧姆），可供选择的器材有：电阻箱 $R$（最大阻值 $99999.9\UO$），滑动变阻器 $R_{1}$（最大阻值 $50\UO$），滑动变阻器 $R_{2}$（最大阻值 $5\UkO$），直流电源 $E$（电动势 $3\UV$），开关 1 个，导线若干。

实验步骤如下：
\begin{steps}
	\item
	按电路原理图连接线路；
	\item
	将电阻箱阻值调节为 0，将滑动变阻器的滑片移到与图中最左端所对应的位置，闭合开关 S；
	\item
	调节滑动变阻器，使电压表满偏；
	\item
	保持滑动变阻器滑片的位置不变，调节电阻箱阻值，使电压表的示数为 $2.00\UV$，记下电阻箱的阻值。
\end{steps}
回答下列问题：
\begin{enumerate}
	\item
	实验中应选择滑动变阻器 \tkanswer[1.5cm]{}（填“$R_{1}$”或“$R_{2}$”）。
	\item
	根据所示电路将图\ref{2016全国理综Ⅱ23b}中实物图连线。
	\item
	实验步骤④中记录的电阻箱阻值为 $630.0\UO$，若认为调节电阻箱时滑动变阻器上的分压不变，计算可得电压表的内阻为 \tkanswer[2cm]{} $\UO$（结果保留到个位）。
	\item
	如果此电压表是由一个表头和电阻串联构成的，可推断该表头的满刻度电流为 \tkanswer[1.5cm]{}（填正确答案标号）。

	\fourchoices[answer=D]
	{$100\UuA$}
	{$250\UuA$}
	{$500\UuA$}
	{$1\UmA$}
\end{enumerate}
\twopicture[labelA=2016全国理综Ⅱ23a, labelB=2016全国理综Ⅱ23b, h=3.4cm, align=right]
{figs/23a.png}
{figs/23b.png}
%% answer: 
\jdanswer{
\begin{enumerate}
	\item
	$R_{1}$

	\item
	如图

	\item
	$2520\UO$

	\item
	D

\end{enumerate}
}
%% memo:
\memoanswer{
	（1）滑动变阻器使用的是分压式接法，为了调节效果好，选择最大阻值小的滑动变阻器 $R_{1}$。

	（2）由电路图连接实物图，注意滑动变阻器的分压式接法，电压表与电阻箱串联，实物图如图所示。

	\onepicture[w=4.5cm]{figs/23_an.png}

	（3）电阻箱阻值为零时，电压表示数为 $2.5\UV$，认为滑动变阻器上的分压不变，即为 $2.5\UV$，所以当电压表示数为 $2\UV$ 时，电阻箱上分得电压为 $0.5\UV$，由串联电路电压比等于电阻比，故电压表的内阻为电阻箱阻值的 4 倍，则 $R_{V}=4\times630\UO=2520\UO$。

	（4）由欧姆定律可得 $I_{g}=\frac{U_{g}}{R_{V}}=\frac{2.5}{2520}\UA\approx1\UmA$。
}

\item
%% number: 24
%% paperName: 2016年全国理综Ⅱ第 24 题 12 分
%% typeId: 6
%% type: 计算题
%% chapter: 电磁感应
%% point: 电磁感应受力分析
%% method: 无
%% score: 12
%% degree: 550
%% duplicateId: 0
%% body:
如图，水平面（纸面）内间距为 $l$ 的平行金属导轨间接一电阻，质量为 $m$、长度为 $l$ 的金属杆置于导轨上。$t=0$ 时，金属杆在水平向右、大小为 $F$ 的恒定拉力作用下由静止开始运动。$t_{0}$ 时刻，金属杆进入磁感应强度大小为 $B$，方向垂直于纸面向里的匀强磁场区域，且在磁场中恰好能保持匀速运动。杆与导轨的电阻均忽略不计，两者始终保持垂直且接触良好，两者之间的动摩擦因数为 $\mu$。重力加速度大小为 $g$。求：
\begin{enumerate}
	\item
	金属杆在磁场中运动时产生的电动势的大小；
	\item
	电阻的阻值。
\end{enumerate}
\onepicture[w=5.5cm, align=right]{figs/24.png}
%% answer: 
\jdanswer{
\begin{enumerate}
	\item
	$E=\frac{Blt_{0}}{m}(F-\mu mg)$

	\item
	$R=\frac{B^{2}l^{2}t_{0}}{m}$

\end{enumerate}
}
%% memo:
\memoanswer{
	（1）设金属杆进入磁场前的加速度大小为 $a$，由牛顿第二定律得
	$F-\mu mg=ma$ \ccnum{①}，

	设金属杆到达磁场左边界时的速度为 $v$，由运动学公式得
	$v=at_{0}$ \ccnum{②}，

	当金属杆以速度 $v$ 在磁场中运动时，由法拉第电磁感应定律可知，杆中的电动势为
	$E=Blv$ \ccnum{③}，

	联立\ccnum{①}\ccnum{②}\ccnum{③}式可得
	$E=Blt_{0}\left(\frac{F}{m}-\mu g\right)$ \ccnum{④}。

	（2）设金属杆在磁场区域中匀速运动时，金属杆中的电流为 $I$，由欧姆定律得
	$I=\frac{E}{R}$ \ccnum{⑤}，

	式中 $R$ 为电阻的阻值，金属杆所受的安培力为
	$f=BIl$ \ccnum{⑥}，

	因金属杆做匀速运动，由牛顿第二定律得
	$F-\mu mg-f=0$ \ccnum{⑦}，

	联立\ccnum{④}\ccnum{⑤}\ccnum{⑥}\ccnum{⑦}式解得
	$R=\frac{B^{2}l^{2}t_{0}}{m}$ \ccnum{⑧}。
}

\item
%% number: 25
%% paperName: 2016年全国理综Ⅱ第 25 题 20 分
%% typeId: 6
%% type: 计算题
%% chapter: 曲线运动
%% point: 竖直平面圆周运动
%% method: 无
%% score: 20
%% degree: 450
%% duplicateId: 0
%% body:
轻质弹簧原长为 $2l$，将弹簧竖直放置在地面上，在其顶端将一质量为 $5m$ 的物体由静止释放，当弹簧被压缩到最短时，弹簧长度为 $l$，现将该弹簧水平放置，一端固定在 A 点，另一端与物块 P 接触但不连接。AB 是长度为 $5l$ 的水平轨道，B 端与半径 $l$ 的光滑半圆轨道 BCD 相切，半圆的直径 RD 竖直，如图所示，物块 P 与 AB 间的动摩擦因数 $\mu=0.5$。用外力推动物块 P，将弹簧压缩至长度 $l$，然后放开，P 开始沿轨道运动，重力加速度大小为 $g$。
\begin{enumerate}
	\item
	若 P 的质量为 $m$，求 P 到达 B 点时的速度的大小，以及它离开圆轨道后落回到 AB 上的位置与 B 点之间的距离；
	\item
	若 P 能滑上圆轨道，且仍能沿圆轨道滑下，求 P 的质量的取值范围。
\end{enumerate}
\onepicture[w=5.5cm, align=right]{figs/25.png}
%% answer: 
\jdanswer{
\begin{enumerate}
	\item
	$\sqrt{6gl}$，$2\sqrt{2}l$

	\item
	$\frac{5}{3}m\leq m'\leq\frac{5}{2}m$

\end{enumerate}
}
%% memo:
\memoanswer{
	（1）依题意，当弹簧竖直放置，长度被压缩至 $l$ 时，质量为 $5m$ 的物体的动能为零，其重力势能转化为弹簧的弹性势能。由机械能守恒定律，弹簧长度为 $l$ 时的弹性势能为
	$E_{p}=5mgl$ \ccnum{①}，

	设 P 的质量为 $M$，到达 B 点时的速度大小为 $v_{B}$，由能量守恒得
	$E_{p}=\frac{1}{2}Mv_{B}^{2}+\mu Mg\cdot4l$ \ccnum{②}，

	联立\ccnum{①}\ccnum{②}式，取 $M=m$ 并代入题给数据得
	$v_{B}=\sqrt{6gl}$ \ccnum{③}，

	若 P 能沿圆轨道运动到 D 点，其到达 D 点时的向心力不能小于重力，即 P 此时的速度大小 $v$ 应该满足
	$\frac{mv^{2}}{l}-mg\geq0$ \ccnum{④}，

	设 P 滑到 D 点时的速度为 $v_{D}$，由机械能守恒得
	$\frac{1}{2}mv_{B}^{2}=\frac{1}{2}mv_{D}^{2}+mg\cdot2l$ \ccnum{⑤}，

	联立\ccnum{③}\ccnum{⑤}式得
	$v_{D}=\sqrt{2gl}$ \ccnum{⑥}，

	$v_{D}$ 满足\ccnum{④}式要求，故 P 能运动到 D 点，并从 D 点以速度 $v_{D}$ 水平射出。

	设 P 落到轨道 AB 所需的时间为 $t$，由运动学公式得
	$2l=\frac{1}{2}gt^{2}$ \ccnum{⑦}，

	P 落回到 AB 上的位置与 B 点之间的距离为
	$s=v_{D}t$ \ccnum{⑧}，

	联立\ccnum{⑥}\ccnum{⑦}\ccnum{⑧}式得 $s=2\sqrt{2}l$ \ccnum{⑨}。

	（2）为使 P 能滑上圆轨道，它到达 B 点时的速度不能小于零，由\ccnum{①}\ccnum{②}式可知
	$5mgl>\mu Mg\cdot4l$ \ccnum{⑩}，

	要使 P 仍能沿圆轨道滑回，P 在圆轨道的上升高度不能超过半圆轨道的中点 C，由机械能守恒定律得
	$\frac{1}{2}Mv_{B}^{2}\leq Mgl$ \ccnum{⑪}，

	联立\ccnum{①}\ccnum{②}\ccnum{⑩}\ccnum{⑪}式得 $\frac{5}{3}m\leq M\leq\frac{5}{2}m$ \ccnum{⑫}。
}

\item
%% number: 33
%% paperName: 2016年全国理综Ⅱ第 33 题 10 分
%% typeId: 6
%% type: 计算题
%% chapter: 气体性质
%% point: 玻意耳定律
%% method: 无
%% score: 10
%% degree: 550
%% duplicateId: 0
%% body:
【物理——选修 3-3】（15 分）

（1）（5 分）一定量的理想气体从状态 a 开始，经历等温或等压过程 ab、bc、cd、da 回到原状态，其 $p-T$ 图像如图所示．其中对角线 ac 的延长线过原点 O。下列判断正确的是\xzanswer{ABE}。（填正确答案标号。选对 1 个得 2 分，选对 2 个得 4 分，选对 3 个得 5 分。每选错 1 个扣 3 分，最低得分为 0 分）

\onepicture[w=4cm]{figs/33.png}

\fivechoices[answer=ABE]
{气体在 a、c 两状态的体积相等}
{气体在状态 a 时的内能大于它在状态 c 时的内能}
{在过程 cd 中气体向外界放出的热量大于外界对气体做的功}
{在过程 da 中气体从外界吸收的热量小于气体对外界做的功}
{在过程 bc 中外界对气体做的功等于在过程 da 中气体对外界做的功}

（2）（10 分）一氧气瓶的容积为 $0.08\Umc$，开始时瓶中氧气的压强为 20 个大气压．某实验室每天消耗 1 个大气压的氧气 $0.36\Umc$。当氧气瓶中的压强降低到 2 个大气压时，需重新充气。若氧气的温度保持不变，求这瓶氧气重新充气前可供该实验室使用多少天。
%% answer: 
\jdanswer{
4 天
}
%% memo:
\memoanswer{
	（1）由理想气体状态方程得 $\frac{pV}{T}=C$，可知 a、c 两状态 $V$ 相等，则 a、c 两状态体积相等，故 A 正确；理想气体在状态 a 时比在状态 c 时温度高，内能大，故 B 正确；由热力学第一定律得 $\Delta U=W+Q$，过程 cd 为等温变化，温度不变，气体的内能不变，$\Delta U=0$，即气体向外界放出的热量等于外界对气体做的功，故 C 错误；过程 da 为等压变化，气体升温膨胀，内能增加，$\Delta U>0$，同时还对外界做功，$W<0$，则由 $\Delta U=W+Q$ 可知 $Q>0$，即必须从外界吸热，且从外界吸收的热量大于气体对外界做的功，故 D 错误；由于 a、c 两状态体积相等，由热力学第一定律 $\Delta U=W+Q$，在过程 bc 中外界对气体做的功等于过程 da 中气体对外界做的功，故 E 正确。

	（2）设氧气开始时的压强为 $p_{1}$，体积为 $V_{1}$，压强变为 $p_{2}$（2 个大气压）时，体积为 $V_{2}$，由玻意耳定律得
	$p_{1}V_{1}=p_{2}V_{2}$ \ccnum{①}，

	重新充气前，用去的氧气在 $p_{2}$ 压强下的体积为
	$V_{3}=V_{2}-V_{1}$ \ccnum{②}，

	设用去的氧气在 $p_{0}$（1 个大气压）压强下的体积为 $V_{0}$，则有
	$p_{2}V_{3}=p_{0}V_{0}$ \ccnum{③}，

	设实验室每天用去的氧气在 $p_{0}$ 下的体积为 $\Delta V$，则氧气可用的天数为
	$N=\frac{V_{0}}{\Delta V}$ \ccnum{④}，

	联立\ccnum{①}\ccnum{②}\ccnum{③}\ccnum{④}式，并代入数据得 $N=4$（天）\ccnum{⑤}。
}

\item
%% number: 34
%% paperName: 2016年全国理综Ⅱ第 34 题 10 分
%% typeId: 6
%% type: 计算题
%% chapter: 机械波
%% point: 波长频率波速
%% method: 无
%% score: 10
%% degree: 450
%% duplicateId: 0
%% body:
【物理——选修 3-4】（15 分）

（1）（5 分）关于电磁波，下列说法正确的是\xzanswer{ABC}。（填正确答案标号。选对 1 个得 2 分，选对 2 个得 4 分，选对 3 个得 5 分。每选错 1 个扣 3 分，最低得分为 0 分）

\fivechoices[answer=ABC]
{电磁波在真空中的传播速度与电磁波的频率无关}
{周期性变化的电场和磁场可以相互激发，形成电磁波}
{电磁波在真空中自由传播时，其传播方向与电场强度、磁感应强度均垂直}
{利用电磁波传递信号可以实现无线通信，但电磁波不能通过电缆、光缆传输}
{电磁波可以由电磁振荡产生，若波源的电磁振荡停止，空间的电磁波随即消失}

（2）（10 分）一列简谐横波在介质中沿 $x$ 轴正向传播，波长不小于 $10\Ucm$。O 和 A 是介质中平衡位置分别位于 $x=0$ 和 $x=5\Ucm$ 处的两个质点。$t=0$ 时开始观测，此时质点 O 的位移为 $y=4\Ucm$，质点 A 处于波峰位置；$t=\frac{1}{3}\Us$ 时，质点 O 第一次回到平衡位置，$t=1\Us$ 时，质点 A 第一次回到平衡位置。求：
\begin{enumerate}
	\item
	简谐波的周期、波速和波长；
	\item
	质点 O 的位移随时间变化的关系式。
\end{enumerate}
%% answer: 
\jdanswer{
\begin{enumerate}
	\item
	$T=4\Us$，$v=7.5\Ucms$，$\lambda=30\Ucm$

	\item
	$y=0.08\sin\left(\frac{\pi}{2}t+\frac{5}{6}\pi\right)$ 或者 $y=0.08\cos\left(\frac{\pi}{2}t+\frac{1}{3}\pi\right)$

\end{enumerate}
}
%% memo:
\memoanswer{
	（1）电磁波在真空中的传播速度为光速 $c$，与频率无关，故 A 正确；周期性变化的电场和磁场可相互激发形成电磁波，故 B 正确；电磁波是横波，在真空中自由传播时传播方向与电场方向、磁场方向垂直，故 C 正确；电磁波可以在介质中传播，所以可以通过光缆进行有线传输，也可以不需要介质进行传输，即无线传输，故 D 错误；电磁振荡停止，已产生的电磁波仍能继续传播而独立存在，故 E 错误。

	（2）（i）设振动周期为 $T$，由于质点 A 在 0 到 $1\Us$ 内由最大位移处第一次回到平衡位置，经历的是 $\frac{1}{4}$ 个周期，由此可知 $T=4\Us$ \ccnum{①}；由于质点 O 与 A 的距离为 $5\Ucm$ 小于半个波长，且波沿着 $x$ 轴正向传播，O 在 $t=\frac{1}{3}\Us$ 时回到平衡位置，而 A 在 $t=1\Us$ 时回到平衡位置，时间 $\frac{2}{3}\Us$。两质点平衡位置的距离除以传播时间，可得到波的传播速度 $v=7.5\Ucms$ \ccnum{②}；利用波长、波速和周期的关系得，简谐波的波长 $\lambda=30\Ucm$ \ccnum{③}。

	（ii）设质点 O 的位移随时间变化的关系为
	$y=A\cos\left(\frac{2\pi t}{T}+\varphi_{0}\right)$ \ccnum{④}，

	将\ccnum{①}式及题给的条件代入上式得
	$4\Ucm=A\cos\varphi_{0}$，$0=A\cos\left(\frac{\pi}{6}+\varphi_{0}\right)$ \ccnum{⑤}，

	解得 $\varphi_{0}=\frac{\pi}{3}$，$A=8\Ucm$ \ccnum{⑥}，

	质点 O 的位移随时间变化的关系式为
	$y=0.08\cos\left(\frac{\pi t}{2}+\frac{\pi}{3}\right)$（国际单位制），

	或 $y=0.08\sin\left(\frac{\pi t}{2}+\frac{5\pi}{6}\right)$（国际单位制）。
}

\item
%% number: 35
%% paperName: 2016年全国理综Ⅱ第 35 题 10 分
%% typeId: 6
%% type: 计算题
%% chapter: 运动和力
%% point: 动量和能量
%% method: 无
%% score: 10
%% degree: 450
%% duplicateId: 0
%% body:
【物理——选修 3-5】（15 分）

（1）（5 分）在下列描述核过程的方程中，属于 $\alpha$ 衰变的是\_\_\_\_\_\_，属于 $\beta$ 衰变的是\_\_\_\_\_\_，属于裂变的是\_\_\_\_\_\_\_\_，属于聚变的是\_\_\_\_\_\_\_\_\_\_\_\xzanswer{CABEF}。

\sixchoices[answer=CABEF]
{\ce{^{14}_{6}C -> ^{14}_{7}N + ^{0}_{-1}e}}
{\ce{^{32}_{15}P -> ^{32}_{16}S + ^{0}_{-1}e}}
{\ce{^{238}_{92}U -> ^{234}_{90}Th + ^{4}_{2}He}}
{\ce{^{14}_{7}N + ^{4}_{2}He -> ^{17}_{8}O + ^{1}_{1}H}}
{\ce{^{235}_{92}U + ^{1}_{0}n -> ^{140}_{54}Xe + ^{94}_{38}Sr + 2^{1}_{0}n}}
{\ce{^{3}_{1}H + ^{2}_{1}H -> ^{4}_{2}He + ^{1}_{0}n}}

（2）（10 分）如图，光滑冰面上静止放置一表面光滑的斜面体，斜面体右侧一蹲在滑板上的小孩和其前面的冰块均静止于冰面上。某时刻小孩将冰块以相对冰面 $3\Ums$ 的速度向斜面体推出，冰块平滑地滑上斜面体，在斜面体上上升的最大高度为 $h=0.3\Um$（$h$ 小于斜面体的高度）。已知小孩与滑板的总质量为 $m_{1}=30\Ukg$，冰块的质量为 $m_{2}=10\Ukg$，小孩与滑板始终无相对运动。取重力加速度的大小 $g=10\Umsq$。
\begin{enumerate}
	\item
	求斜面体的质量；
	\item
	通过计算判断，冰块与斜面体分离后能否追上小孩？
\end{enumerate}
\onepicture[w=7cm, align=right]{figs/35.png}
%% answer: 
\jdanswer{
\begin{enumerate}
	\item
	$M=20\Ukg$

	\item
	追不上

\end{enumerate}
}
%% memo:
\memoanswer{
	（1）A、B 是 $\beta$ 衰变；C 是 $\alpha$ 衰变；D 是人工转变；E 是裂变；F 是聚变。

	（2）（i）规定向右为速度正方向。冰块在斜面体上运动到最大高度时两者达到共同速度，设此共同速度为 $v$，斜面体的质量为 $m_{3}$，由水平方向动量守恒和机械能守恒得
	$m_{2}v_{20}=(m_{2}+m_{3})v$ \ccnum{①}，

	$\frac{1}{2}m_{2}v_{20}^{2}=\frac{1}{2}(m_{2}+m_{3})v^{2}+m_{2}gh$ \ccnum{②}，

	式中 $v_{20}=-3\Ums$ 为冰块推出时的速度，

	联立\ccnum{①}\ccnum{②}式并代入题给数据得
	$m_{3}=20\Ukg$ \ccnum{③}。

	（ii）设小孩推出冰块后的速度为 $v_{1}$，由动量守恒定律得
	$m_{1}v_{1}+m_{2}v_{20}=0$ \ccnum{④}，

	代入数据得 $v_{1}=1\Ums$ \ccnum{⑤}，

	设冰块与斜面体分离后的速度分别为 $v_{2}$ 和 $v_{3}$，由动量守恒和机械能守恒得
	$m_{2}v_{20}=m_{2}v_{2}+m_{3}v_{3}$ \ccnum{⑥}，

	$\frac{1}{2}m_{2}v_{20}^{2}=\frac{1}{2}m_{2}v_{2}^{2}+\frac{1}{2}m_{3}v_{3}^{2}$ \ccnum{⑦}，

	联立\ccnum{③}\ccnum{⑥}\ccnum{⑦}式并代入数据得 $v_{2}=1\Ums$ \ccnum{⑧}，

	由于冰块与斜面体分离后的速度与小孩推出冰块后的速度相同且处在后方，故冰块不能追上小孩。
}

\end{enumerate}

\end{document}
